\chapter{Breast Augmentation}

\section*{Introduction \& Clinical Indications}
Breast augmentation is an elective cosmetic surgical procedure designed to increase the size and improve the shape of the breasts, primarily for aesthetic enhancement. This operation involves the surgical placement of medical implants, typically filled with saline or silicone gel, or less commonly, the transfer of autologous fat. The procedure is performed to address concerns such as naturally small breasts, loss of breast volume following pregnancy, breastfeeding, or significant weight loss, and to correct congenital or acquired breast asymmetry. As one of the most frequently performed cosmetic surgeries globally, its clinical significance lies in its profound impact on a patient's body image, self-confidence, and overall psychological well-being, rather than treating a specific disease state.

\begin{itemize}
  \item Augmentation mammoplasty
  \item Breast implant surgery
  \item Breast enhancement surgery
  \item Mammoplasty augmentation
  \item Implant-based breast enlargement
  \item Fat transfer breast augmentation (when autologous fat is used)
\end{itemize}

The fundamental principle of breast augmentation is to increase the volume and projection of the breast by introducing a prosthetic device or autologous tissue into a surgically created pocket. For implant-based augmentation, the implant acts as a permanent volume expander, placed either directly behind the glandular tissue (subglandular or prepectoral), beneath the pectoralis major muscle (submuscular or retropectoral), or in a dual-plane configuration that partially covers the implant with muscle. The rationale for choosing an implant plane depends on the patient's existing breast tissue, desired aesthetic outcome, and surgeon's preference, aiming to optimize implant coverage and minimize palpability.

The primary technical goals include achieving a natural-looking and symmetric result, ensuring appropriate implant positioning relative to the inframammary fold and nipple, and minimizing the risk of complications such as capsular contracture or implant malposition. Autologous fat transfer, while less common for significant volume increase, operates on the principle of harvesting viable adipose tissue from donor sites (e.g., abdomen, flanks) and grafting it into the breast, offering a natural tissue solution without synthetic materials. Both approaches aim to enhance the patient's breast contour and proportion in harmony with their overall physique.

The quest for breast augmentation dates back to the late 19th and early 20th centuries, with early attempts involving the injection of various foreign substances such as paraffin, beeswax, and even glass balls, which often led to severe complications and granuloma formation. The modern era of breast augmentation began in 1962 when American plastic surgeons Thomas Cronin and Frank Gerow developed the first silicone gel breast implant, marking a significant milestone in reconstructive and cosmetic surgery. Their pioneering work at the Dow Corning Corporation revolutionized the field, offering a safer and more predictable method for breast enlargement.

Over the subsequent decades, breast implant technology underwent continuous refinement. Saline-filled implants were introduced as an alternative, and silicone gel implants evolved with improved shell designs, including textured surfaces and more cohesive gel formulations to reduce rupture rates and gel bleed. In 1992, the U.S. Food and Drug Administration (FDA) placed a moratorium on silicone gel implants for cosmetic use due to safety concerns, leading to a surge in saline implant use. Following extensive research and clinical trials, silicone gel implants were reintroduced for cosmetic augmentation in 2006, with updated safety profiles and recommendations for long-term surveillance. Today, ongoing research focuses on implant surface technologies to minimize capsular contracture and the rare risk of implant-associated anaplastic large cell lymphoma (BIA-ALCL).

Breast augmentation is primarily indicated for elective cosmetic enhancement, but also serves reconstructive purposes.

\begin{itemize}
  \item To increase breast size in women who perceive their breasts as too small or desire a fuller décolletage.
  \item To restore breast volume and projection lost due to pregnancy, breastfeeding, or significant weight fluctuations.
  \item To correct congenital breast asymmetry, where one breast is naturally smaller or differently shaped than the other.
  \item To address congenital breast hypoplasia or amastia, including conditions like Poland syndrome or tuberous breast deformity.
  \item To improve overall breast shape and contour for enhanced body proportion and aesthetic appeal.
  \item As a component of breast reconstruction following lumpectomy or partial mastectomy, to achieve symmetry with the contralateral breast.
  \item To enhance self-image and confidence in individuals with body dysmorphia or dissatisfaction related to breast size.
\end{itemize}

Breast augmentation is predominantly a primary, elective cosmetic procedure, initiated by the patient's desire for aesthetic improvement. It is not typically a medically necessary intervention for disease treatment. In this primary role, it aims to enhance body contour, restore lost volume, or correct inherent size discrepancies.

However, breast augmentation can also serve a secondary or reconstructive role in specific clinical scenarios. For instance, it may be performed as part of a staged breast reconstruction after mastectomy, often in conjunction with tissue expanders, or to achieve symmetry following unilateral breast cancer surgery (e.g., lumpectomy or partial mastectomy) where the remaining breast tissue is insufficient. In cases of congenital breast deformities such as tuberous breasts or Poland syndrome, augmentation is considered reconstructive, aiming to normalize breast shape and volume. The distinction between primary cosmetic and secondary reconstructive positioning influences insurance coverage, patient counseling, and the overall surgical planning.

\section*{Relevant Surgical Anatomy}
A thorough understanding of breast anatomy is paramount for safe and effective breast augmentation. The breast itself is composed of glandular tissue, fibrous connective tissue, and adipose tissue, all encased within the superficial fascia and overlying the pectoralis major muscle. Key anatomical landmarks include the nipple-areola complex, the inframammary fold (IMF), and the lateral breast border. The IMF is a crucial landmark for implant placement, defining the lower border of the breast mound.

The pectoralis major muscle forms the posterior boundary of the breast in the upper two-thirds, originating from the clavicle, sternum, and costal cartilages, and inserting into the humerus. The pectoralis minor muscle lies deep to the pectoralis major. The arterial supply to the breast is rich, primarily from the internal mammary artery (medially), the lateral thoracic artery (laterally), and intercostal perforators. Venous drainage largely parallels the arterial supply. Sensory innervation to the breast and nipple-areola complex is provided by the anterior and lateral cutaneous branches of the second through sixth intercostal nerves. The long thoracic nerve and thoracodorsal nerve, which innervate the serratus anterior and latissimus dorsi muscles respectively, are important structures to identify and protect during submuscular dissection, particularly in the axillary region. Lymphatic drainage is predominantly to the axillary lymph nodes, with some drainage to internal mammary nodes.

\section*{Preoperative Workup \& Preparation}
A comprehensive preoperative workup is essential to ensure patient safety and optimize surgical outcomes. This begins with a detailed consultation where the patient's aesthetic goals, medical history, current medications, and allergies are thoroughly reviewed. A physical examination of the breasts and chest wall is performed to assess existing breast size, shape, skin quality, and tissue elasticity, which guides implant selection and surgical planning.

\begin{itemize}
  \item \textbf{Medical Evaluation:} A general medical evaluation, including a review of cardiovascular and pulmonary health, is conducted to ensure fitness for anesthesia and surgery. This may involve blood tests (e.g., complete blood count, coagulation profile), an electrocardiogram (ECG), or chest X-ray, particularly for older patients or those with comorbidities.
  \item \textbf{Imaging:} Baseline mammography or ultrasound may be recommended, especially for women over 40 or those with a family history of breast cancer, to screen for any underlying breast pathology and establish a baseline for future breast cancer screening.
  \item \textbf{Medication Adjustment:} Patients are advised to discontinue medications that increase bleeding risk, such as aspirin, NSAIDs, and certain herbal supplements (e.g., ginkgo biloba, vitamin E), typically two weeks prior to surgery. Smoking cessation is strongly encouraged for at least 4-6 weeks preoperatively to improve wound healing and reduce complications.
  \item \textbf{NPO Status:} Patients are instructed to fast (NPO - nil per os) for a specified period, usually 6-8 hours, before surgery to prevent aspiration during anesthesia.
  \item \textbf{Prophylactic Antibiotics:} Intravenous prophylactic antibiotics are administered shortly before the surgical incision to minimize the risk of surgical site infection.
  \item \textbf{Informed Consent:} A detailed discussion of the procedure, expected outcomes, potential risks, and the need for future implant surveillance is conducted, and informed consent is obtained.
\end{itemize}

Careful patient selection is crucial for successful breast augmentation.

\textbf{Situations requiring treatment, optimization, or an alternative plan:}
\begin{itemize}
  \item Active breast infection or inflammation.
  \item Untreated or active breast cancer.
  \item Known allergy to implant materials (e.g., silicone).
  \item Pregnancy or active breastfeeding.
  \item Unrealistic patient expectations or severe body dysmorphic disorder.
  \item Significant systemic illness that precludes safe anesthesia or surgery (e.g., severe uncontrolled heart disease, uncontrolled diabetes).
\end{itemize}

\textbf{Relative Contraindications \& Precautions:}
\begin{itemize}
  \item Poorly controlled chronic medical conditions (e.g., diabetes, autoimmune diseases) that may impair healing or increase infection risk.
  \item A history of abnormal scarring, such as keloids or hypertrophic scars.
  \item Significant ptosis (sagging) of the breasts, which may require a mastopexy (breast lift) in conjunction with or instead of augmentation.
  \item Connective tissue disorders or conditions affecting wound healing.
  \item Smoking, which significantly increases the risk of complications, including wound healing issues and capsular contracture.
  \item Patients with a strong family history of breast cancer, requiring careful discussion about implant effects on mammographic screening.
  \item The need for future implant surveillance and potential revision surgery over the patient's lifetime.
\end{itemize}

\section*{Surgical Technique \& Procedure}
Breast augmentation is typically performed in an accredited surgical facility under general anesthesia, though local anesthesia with sedation may be used in select cases. The procedure usually lasts one to two hours.

\begin{itemize}
  \item \textbf{Anesthesia and Positioning:} The patient is placed in a supine position. Preoperative markings are meticulously drawn while the patient is upright, outlining the inframammary fold, nipple-areola complex, and desired implant pocket boundaries. Prophylactic intravenous antibiotics are administered prior to incision.
  \item \textbf{Incision:} The surgeon makes an incision in a discreet location to minimize visible scarring. Common incision sites include:
    \begin{itemize}
      \item \textbf{Inframammary Fold (IMF):} The most common approach, offering direct access and excellent visualization, with the scar hidden in the natural breast crease.
      \item \textbf{Periareolar:} An incision made around the lower half of the areola, camouflaged by the change in skin pigmentation.
      \item \textbf{Transaxillary:} An incision made in the armpit, allowing for remote placement of the implant without breast incisions, often requiring endoscopic assistance.
    \end{itemize}
  \item \textbf{Pocket Creation:} Through the incision, a surgical pocket is carefully dissected to accommodate the implant. The plane of dissection can be:
    \begin{itemize}
      \item \textbf{Subglandular (Prepectoral):} Directly behind the breast glandular tissue, on top of the pectoralis major muscle.
      \item \textbf{Submuscular (Retropectoral):} Beneath the pectoralis major muscle, providing more tissue coverage over the implant.
      \item \textbf{Dual-Plane:} A hybrid approach where the upper portion of the implant is covered by the pectoralis muscle, and the lower portion lies subglandular.
    \end{itemize}
  \item \textbf{Implant Insertion:} The chosen implant (saline or silicone gel) is carefully inserted into the created pocket. Saline implants are typically inserted empty and then filled to the desired volume with sterile saline solution. Silicone gel implants are pre-filled. The surgeon ensures proper positioning and symmetry.
  \item \textbf{Hemostasis and Irrigation:} The pocket is meticulously inspected for bleeding, and hemostasis is achieved. The pocket may be irrigated with an antibiotic solution to reduce infection risk.
  \item \textbf{Closure:} The incisions are closed in layers using absorbable sutures for the deeper tissues and fine sutures or surgical glue for the skin. Drains are rarely used in primary augmentation but may be considered in revision cases or if significant bleeding is encountered.
  \item \textbf{Dressing:} A supportive surgical bra or dressing is applied to provide compression and support during the initial healing phase.
\end{itemize}

\section*{Complications \& Risks}
While generally safe, breast augmentation carries potential complications, both general surgical and procedure-specific.

\textbf{General Surgical Complications:}
\begin{itemize}
  \item \textbf{Anesthesia Risks:} Reactions to anesthesia, including nausea, vomiting, allergic reactions, or, rarely, more severe cardiovascular or respiratory events.
  \item \textbf{Bleeding/Hematoma:} Accumulation of blood within the surgical site, potentially requiring reoperation for evacuation.
  \item \textbf{Infection:} Bacterial infection of the surgical site or implant, which may necessitate antibiotic treatment or implant removal.
  \item \textbf{Seroma:} Accumulation of clear fluid around the implant, sometimes requiring drainage.
  \item \textbf{Poor Wound Healing/Scarring:} Delayed healing, hypertrophic scars, or keloids, particularly in predisposed individuals.
\end{itemize}

\textbf{Procedure-Specific Complications:}
\begin{itemize}
  \item \textbf{Capsular Contracture:} The most common long-term complication, where the scar tissue (capsule) around the implant hardens and tightens, leading to breast firmness, pain, distortion, and asymmetry. It is graded from I (soft) to IV (hard, painful, distorted).
  \item \textbf{Implant Rupture or Leakage:} Failure of the implant shell, leading to leakage of saline or silicone gel. Saline rupture is usually obvious due to breast deflation; silicone rupture can be "silent" and may require imaging (MRI) for diagnosis.
  \item \textbf{Changes in Nipple or Breast Sensation:} Temporary or permanent numbness, hypersensitivity, or altered sensation in the nipples or breast skin.
  \item \textbf{Asymmetry or Unsatisfactory Cosmetic Result:} Differences in breast size, shape, or position, or a result that does not meet patient expectations, potentially requiring revision surgery.
  \item \textbf{Implant Malposition:} Displacement of the implant, such as superior displacement (high-riding), inferior displacement (bottoming out), or lateral displacement.
  \item \textbf{Rippling/Wrinkling:} Visible or palpable folds in the implant shell, more common with saline implants or in patients with minimal breast tissue.
  \item \textbf{Breast Implant-Associated Anaplastic Large Cell Lymphoma (BIA-ALCL):} A rare, but serious, T-cell lymphoma associated primarily with textured breast implants. It is typically curable with implant and capsule removal.
  \item \textbf{Breast Pain:} Chronic or intermittent breast pain, which can be challenging to manage.
  \item \textbf{Need for Reoperation:} Implants are not lifetime devices, and revision surgery for complications, aesthetic changes, or implant exchange is common over time.
\end{itemize}

\section*{Postoperative Care \& Recovery}
The postoperative recovery timeline for breast augmentation typically involves several phases, with most patients resuming light activities within a week and full recovery taking several months.

Immediately after surgery, patients are monitored in the Post-Anesthesia Care Unit (PACU) for vital signs, pain control, and recovery from anesthesia. Most patients are discharged home the same day. During the first 24-48 hours, patients will experience swelling, bruising, and moderate pain, which is managed with prescribed oral analgesics. A supportive surgical bra is worn continuously, and strenuous activities, heavy lifting, and arm elevation above the shoulders are strictly avoided. Rest is paramount, and patients are encouraged to walk lightly to promote circulation.

Over the first 1-2 weeks, discomfort gradually subsides, and bruising begins to resolve. Most patients can return to light, non-strenuous work within 5-7 days, depending on their occupation. Sutures, if non-dissolvable, are typically removed at a follow-up visit around 7-14 days post-op. Patients are usually advised to continue wearing the supportive bra for 4-6 weeks. Strenuous exercise, heavy lifting, and impact activities are generally restricted for 4-6 weeks to allow for proper tissue healing and implant stabilization. Long-term recovery involves the gradual softening of the breast tissue and implants, with the final aesthetic result often becoming apparent after 3-6 months as swelling fully resolves and the implants "drop and fluff" into their final position.

During breast augmentation, the surgeon's findings primarily relate to the existing breast tissue characteristics and the creation of the implant pocket. Key observations include the quality and quantity of native breast tissue, the presence of any asymmetry, the integrity of the inframammary fold, and the ease or difficulty of dissecting the chosen implant pocket plane (subglandular, submuscular, or dual-plane). The surgeon meticulously ensures hemostasis and confirms the correct size and positioning of the implants within the created space. Any unexpected findings, such as unusual tissue texture or incidental masses, would be documented and potentially biopsied.

For primary breast augmentation, tissue is generally not resected, so a pathology report is typically not generated unless a pre-existing lesion was identified and biopsied. In cases of revision surgery, however, a portion or the entirety of the periprosthetic capsule (scar tissue surrounding the implant) may be removed (capsulectomy) and sent for pathological examination. This is particularly important if there is suspicion of capsular contracture, infection, or, critically, breast implant-associated anaplastic large cell lymphoma (BIA-ALCL), where the pathology report would confirm the diagnosis and characterize the cellular findings.

A normal and successful postoperative course following breast augmentation is characterized by a predictable progression of healing and resolution of initial symptoms. Patients can expect initial swelling and bruising to gradually diminish over the first few weeks. Pain, which is typically moderate in the first few days, should steadily improve with prescribed oral pain medication and eventually resolve. The breasts will initially feel firm and may appear higher than their final position, but over several weeks to months, the implants will "drop and fluff" into a more natural, softer contour as the tissues relax and swelling subsides.

The surgical incisions will heal, and scars will mature, gradually fading from red to pink to a fine white line over 6-12 months. Nipple and breast sensation may be temporarily altered but often returns to normal or near-normal within several months, though permanent changes are possible. The patient should feel increasingly comfortable and able to resume normal daily activities, with a growing satisfaction in the aesthetic outcome. A successful course is marked by the absence of significant complications such as infection, hematoma, or early capsular contracture, and the achievement of symmetric, natural-looking breasts that align with the patient's preoperative goals.

Diligent postoperative care and regular follow-up are crucial for optimal healing and long-term implant surveillance.

\begin{itemize}
  \item \textbf{Follow-up Visits:} Initial follow-up appointments are typically scheduled within the first week to check wound healing, remove any non-dissolvable sutures, and assess for early complications. Subsequent visits are usually at 1 month, 3 months, 6 months, and annually thereafter.
  \item \textbf{Wound Care:} Patients are instructed on how to care for their incisions, keeping them clean and dry. Scar management techniques, such as silicone sheeting or massage, may be recommended once incisions are fully healed.
  \item \textbf{Activity Restrictions:} Adherence to activity restrictions, including avoiding heavy lifting and strenuous exercise for 4-6 weeks, is vital to prevent complications like hematoma, seroma, or implant displacement.
  \item \textbf{Supportive Bra:} Wearing a supportive surgical bra continuously for several weeks, as advised by the surgeon, helps to stabilize the implants and reduce swelling.
  \item \textbf{Implant Massage:} Some surgeons recommend specific implant massage protocols to help prevent capsular contracture, particularly with smooth implants.
  \item \textbf{Warning Signs:} Patients are educated on warning signs that warrant immediate medical attention, including fever, increasing pain, excessive swelling, redness, pus drainage from incisions, or sudden changes in breast shape or sensation.
  \item \textbf{Long-term Surveillance:} Regular self-breast exams and routine mammographic screening for breast cancer should continue as recommended for the patient's age and risk factors. For silicone gel implants, the FDA recommends MRI screening for silent rupture 5-6 years after implantation and every 2-3 years thereafter, although this recommendation is often debated.
\end{itemize}
Implants have a finite lifespan, and patients should be aware that revision surgery may be necessary in the future due to complications or simply the desire for a change.

\section*{Outcomes \& Prognosis}
Breast augmentation consistently ranks among the most frequently performed cosmetic surgical procedures worldwide. In the United States, hundreds of thousands of procedures are performed annually, making it a cornerstone of aesthetic surgery. The majority of patients are women, typically ranging in age from their early 20s to their 50s. Younger women often seek augmentation for naturally small breasts or to achieve a desired aesthetic, while older women frequently undergo the procedure to restore volume lost due to pregnancy, breastfeeding, or age-related changes.

Silicone gel implants are the most popular choice, accounting for a larger proportion of augmentations compared to saline implants. The incidence of complications varies, with capsular contracture being the most common long-term issue, affecting an estimated 10-15\% of patients over 10 years, though rates vary based on implant type, placement, and patient factors. The overall morbidity of breast augmentation is low when performed by experienced, board-certified plastic surgeons, and serious complications are rare. Mortality rates are exceedingly low, comparable to other elective surgical procedures. The prevalence of breast implant-associated anaplastic large cell lymphoma (BIA-ALCL) is extremely rare, estimated at 1 in 2,200 to 1 in 86,000 for textured implants, with a much lower risk for smooth implants.

The outcomes of breast augmentation are generally highly favorable, with high rates of patient satisfaction reported in numerous studies. Most women achieve their desired aesthetic goals, experiencing improved body image, self-confidence, and overall quality of life. The cosmetic results are typically durable, providing long-lasting enhancement. However, it is crucial for patients to understand that breast implants are not lifetime devices. The prognosis for implant longevity is finite, with a significant proportion of patients requiring at least one revision surgery over their lifetime, often within 10-15 years, due to complications such as capsular contracture, implant rupture, or simply a desire for a change in size or shape.

Prognostic factors for successful outcomes include appropriate patient selection, realistic expectations, meticulous surgical technique, and adherence to postoperative care instructions. While complications like capsular contracture can impact the aesthetic result and may necessitate further surgery, the prognosis for managing these issues is generally good. For the rare but serious complication of BIA-ALCL, early diagnosis and complete surgical removal of the implant and capsule typically lead to an excellent prognosis. Landmark studies and long-term data from organizations like the American Society of Plastic Surgeons consistently demonstrate the safety and efficacy of breast augmentation, affirming its role as a highly successful procedure for aesthetic and reconstructive purposes.

\section*{References}
\begin{enumerate}
  \item American Society of Plastic Surgeons. Breast Augmentation. Available from: \url{https://www.plasticsurgery.org/cosmetic-procedures/breast-augmentation}. Accessed 2024.
  \item Sabiston Textbook of Surgery: The Biological Basis of Modern Surgical Practice. 22nd ed. Elsevier; 2022.
  \item Spear SL, et al. The American Society of Plastic Surgeons 2018 Breast Implant Consensus Statement. Plast Reconstr Surg. 2019 Jan;143(1):1S-10S.
  \item MedlinePlus. Breast augmentation. U.S. National Library of Medicine; 2024. Available from: \url{https://medlineplus.gov/breastaugmentation.html}. Accessed 2024.
\end{enumerate}

\clearpage
